\section{Introduction}
\label{sec:intro}

The Moon is among the best-observed bodies in the solar system. Since 2009 the Lunar Reconnaissance Orbiter (LRO) has returned camera images~\cite{robinson2010lunar}, laser altimetry~\cite{smith2010lunar}, thermal radiometry~\cite{paige2010lunar} and radar~\cite{nozette2010lunar}; Kaguya added multiband imagery and stereo topography~\cite{kato2010kaguya,ohtake2008performance}, GRAIL mapped the gravity field~\cite{zuber2013gravity} and Lunar Prospector the near-surface hydrogen~\cite{binder1998lunar,lawrence2022global}. Yet this archive is not a dataset. Camera images are stored as raw experiment data records, uncalibrated and unprojected, each with its own viewing and illumination geometry. Global maps come at resolutions from 60\,m to 20\,km, in different projections, units and missing-value conventions. Before a model can learn from several instruments at once, their products have to be brought onto the same ground.

Earth observation faced a milder version of this problem and answered it with foundation models pretrained on co-registered multi-sensor archives~\cite{cong2022satmae,fuller2023croma,jakubik2025terramind,tseng2025galileo,astruc2025anysat}. The first lunar foundation models appeared in 2026~\cite{fraccaro2026multimodal,gironamata2026lunarfm,prasad2026moonstone}. All three are trained by reconstructing masked pixels or tokens, resample their inputs to a shared grid, and treat each tile as an independent sample.

We take a different view. The lunar surface is a persistent state: topography, albedo, roughness and composition change on geological time scales. Static maps measure this state directly. A camera image is a rendering of it under a particular sun, and on an airless body the rendering is governed by little other than topography, albedo and illumination. Observations of the same footprint under different suns are therefore views of one state, and a model that must predict one view from another, given the geometry of the target, has to represent the terrain and how light acts on it. This is the structure of a world model~\cite{ha2018recurrent,lecun2022path}: a state estimated from observations and a predictor that rolls the state forward under a control variable. Here the control variable is the acquisition geometry rather than an action, and there is no time axis.

\paragraph{Terminology.} A \emph{footprint} is a fixed 60\,km square of ground; it is the sample. A \emph{modality} is one kind of input, such as a WAC image, elevation or mineralogy; planetary science calls the archived files these come from \emph{products}, and we keep that word for the archive only. A \emph{static map} is a modality that does not depend on the sun, so a footprint has one copy of it (\cref{sec:dataset} calls a stored map a layer, since one map can hold several bands). A \emph{camera observation} is one image of the footprint together with its \emph{acquisition geometry}, the six angles that fix the sun and the viewing direction; a footprint has many. An \emph{observation} is either a static map or a camera observation. A footprint is divided into a $G\times G$ grid of \emph{ground cells}, and every modality is tokenized so that one token covers one ground cell. The \emph{state} is what all observations of a footprint share: one latent vector per ground cell that the model has to estimate.

Learning such a model needs data that the archive does not provide: many observations of the same ground under different suns, each with its geometry, co-registered with the static maps at their native resolutions. SomBench~\cite{patil2026sombench} already does the first step: it calibrates the raw LRO camera images, projects them onto the SLDEM2015 terrain model so that every pixel sits at its true map position, and reprojects the global maps into the same Lunar Transverse Mercator (LTM) zones. Its tiles are cut from one image at a time, however: a tile is a single picture of a place, and two pictures of the same place become two unrelated tiles with different bounds. On these inputs we build a footprint-anchored corpus (\cref{sec:dataset}). Its 10,629 footprints of 60\,km are anchored to the GRAIL gravity grid. The 30 static maps are cut once per footprint at native resolution and shared by all its observations. Each footprint keeps up to 500 Wide Angle Camera (WAC) observations on identical bounds, selected across bins of incidence and sun azimuth, about 52 per covered footprint on average. It is subdivided into 500\,m sub-tiles that carry 1\,m Narrow Angle Camera (NAC) images and 3\,m stereo terrain models. Splits follow LTM zones, and a seam rule ensures that no two footprints, and hence no two splits, share ground.

On this corpus we train \method (\cref{sec:method}), a joint-embedding predictive architecture~\cite{assran2023self,bardes2024revisiting}. Each modality is tokenized by its own stem at its native resolution, so that a token covers the same ground cell in every modality. An encoder reads a partially masked observation of a footprint, without being told its sun. A predictor, given the target modality and its sun, predicts the token grid that an exponential moving average (EMA) of the encoder produces from the actual target observation. Nothing is reconstructed in pixel space.

This kind of model can fail in a way that is easy to miss (\cref{sec:collapse}). Much of a target is unpredictable from the source at the token level, while source and target share information about token position and modality identity. The model can lower its loss by encoding that shared information instead of the ground. We call the result a position code. It is related to the positional collapse that CAPI reported for latent masked image modeling~\cite{darcet2025cluster}, but here it comes with a falling loss, a rising within-image rank and a shrinking VICReg penalty. We explain why, show that two standard regularizers favor it, and give a cross-sample monitor that costs one cosine per target.

Our contributions are:
\begin{itemize}
  \item a footprint-anchored corpus on the SomBench inputs: GRAIL-grid footprints with static maps cut once at native pixels, up to 500 illumination-diverse WAC observations per footprint on identical bounds with their geometry, 500\,m sub-tiles with NAC, stereo terrain and per-sub-tile static-map means, and zone splits whose seam ownership leaves no shared ground (\cref{sec:dataset});
  \item \method, to our knowledge the first world model of the lunar surface. The ground is the sample and the acquisition geometry plays the role of the action: the model estimates a latent state of a footprint from a partially masked observation and, told which modality and which sun, predicts another observation in embedding space. Illumination thus becomes an input the model must explain rather than a nuisance it averages out, and one state serves camera images and instrument maps alike at their native resolutions (\cref{sec:method});
  \item a characterization of position-code collapse in a cross-modal JEPA with continuous targets, an explanation of why the loss, within-tile rank and VICReg penalties cannot see it, and a cross-sample monitor that detected all seven failures we analyzed (\cref{sec:collapse});
  \item experiments on crater detection, shape and albedo from shading and a sun-swap test, which show that frozen features match a masked-token lunar foundation model, gain substantially from a second modality, and respond to the sun the predictor is given (\cref{sec:experiments}).
\end{itemize}
